\documentclass[10pt,a4paper]{amsart}
\usepackage{amsmath,amssymb,amsthm,graphicx}
% Retained minimal preamble for the Stacks Project chapter more-morphisms.tex (More on Morphisms).
% Verbatim excerpt of the chapter preamble: Xy-pic setup (native preamble lines 14-18) and the native
% theorem declarations with the mathematical macros (native preamble lines 154-217). Copied mechanically;
% nothing added, removed or redefined. The wrapper class and the amsmath/amssymb/amsthm/graphicx base
% packages are supplied by the assembler, so the class/externaldocument part of the native preamble is omitted.
% For commutative diagrams we use Xy-pic
\usepackage[all]{xy}

% We use 2cell for 2-commutative diagrams.
\xyoption{2cell}
% Theorem environments.
%
\theoremstyle{plain}
\newtheorem{theorem}[subsection]{Theorem}
\newtheorem{proposition}[subsection]{Proposition}
\newtheorem{lemma}[subsection]{Lemma}

\theoremstyle{definition}
\newtheorem{definition}[subsection]{Definition}
\newtheorem{example}[subsection]{Example}
\newtheorem{exercise}[subsection]{Exercise}
\newtheorem{situation}[subsection]{Situation}

\theoremstyle{remark}
\newtheorem{remark}[subsection]{Remark}
\newtheorem{remarks}[subsection]{Remarks}

\numberwithin{equation}{subsection}

% Macros
%
\def\lim{\mathop{\mathrm{lim}}\nolimits}
\def\colim{\mathop{\mathrm{colim}}\nolimits}
\def\Spec{\mathop{\mathrm{Spec}}}
\def\Hom{\mathop{\mathrm{Hom}}\nolimits}
\def\Ext{\mathop{\mathrm{Ext}}\nolimits}
\def\SheafHom{\mathop{\mathcal{H}\!\mathit{om}}\nolimits}
\def\SheafExt{\mathop{\mathcal{E}\!\mathit{xt}}\nolimits}
\def\Sch{\mathit{Sch}}
\def\Mor{\mathop{\mathrm{Mor}}\nolimits}
\def\Ob{\mathop{\mathrm{Ob}}\nolimits}
\def\Sh{\mathop{\mathit{Sh}}\nolimits}
\def\NL{\mathop{N\!L}\nolimits}
\def\CH{\mathop{\mathrm{CH}}\nolimits}
\def\proetale{{pro\text{-}\acute{e}tale}}
\def\etale{{\acute{e}tale}}
\def\QCoh{\mathit{QCoh}}
\def\Ker{\mathop{\mathrm{Ker}}}
\def\Im{\mathop{\mathrm{Im}}}
\def\Coker{\mathop{\mathrm{Coker}}}
\def\Coim{\mathop{\mathrm{Coim}}}

% Boxtimes
%
\DeclareMathSymbol{\boxtimes}{\mathbin}{AMSa}{"02}

%
% Macros for moduli stacks/spaces
%
\def\QCohstack{\mathcal{QC}\!\mathit{oh}}
\def\Cohstack{\mathcal{C}\!\mathit{oh}}
\def\Spacesstack{\mathcal{S}\!\mathit{paces}}
\def\Quotfunctor{\mathrm{Quot}}
\def\Hilbfunctor{\mathrm{Hilb}}
\def\Curvesstack{\mathcal{C}\!\mathit{urves}}
\def\Polarizedstack{\mathcal{P}\!\mathit{olarized}}
\def\Complexesstack{\mathcal{C}\!\mathit{omplexes}}
% \Pic is the operator that assigns to X its picard group, usage \Pic(X)
% \Picardstack_{X/B} denotes the Picard stack of X over B
% \Picardfunctor_{X/B} denotes the Picard functor of X over B
\def\Pic{\mathop{\mathrm{Pic}}\nolimits}
\def\Picardstack{\mathcal{P}\!\mathit{ic}}
\def\Picardfunctor{\mathrm{Pic}}
\def\Deformationcategory{\mathcal{D}\!\mathit{ef}}
\newenvironment{reference}{}{}
\newenvironment{slogan}{}{}
\newenvironment{history}{}{}
\usepackage{cleveref}
\crefname{originallemma}{original lemma}{original lemmas}
\Crefname{originallemma}{Original lemma}{Original lemmas}
\makeatletter
\newlabel{lemma-isomorphism-approximation}{{lemma-isomorphism-approximation}{source}{}{sourcepointer-1}{}}
\newlabel{lemma-isomorphism-approximation@cref}{{[originallemma][0][]lemma-isomorphism-approximation}{[1][0][]source}{}{}{}}
\makeatother
\crefname{originallemma}{original lemma}{original lemmas}
\Crefname{originallemma}{Original lemma}{Original lemmas}
\makeatletter
\newlabel{lemma-slice-smooth}{{lemma-slice-smooth}{source}{}{sourcepointer-2}{}}
\newlabel{lemma-slice-smooth@cref}{{[originallemma][0][]lemma-slice-smooth}{[1][0][]source}{}{}{}}
\makeatother
\crefname{originallemma}{original lemma}{original lemmas}
\Crefname{originallemma}{Original lemma}{Original lemmas}
\makeatletter
\newlabel{limits-lemma-descend-etale}{{limits-lemma-descend-etale}{source}{}{sourcepointer-3}{}}
\newlabel{limits-lemma-descend-etale@cref}{{[originallemma][0][]limits-lemma-descend-etale}{[1][0][]source}{}{}{}}
\makeatother
\crefname{originallemma}{original lemma}{original lemmas}
\Crefname{originallemma}{Original lemma}{Original lemmas}
\makeatletter
\newlabel{limits-lemma-descend-finite-presentation}{{limits-lemma-descend-finite-presentation}{source}{}{sourcepointer-4}{}}
\newlabel{limits-lemma-descend-finite-presentation@cref}{{[originallemma][0][]limits-lemma-descend-finite-presentation}{[1][0][]source}{}{}{}}
\makeatother
\crefname{originallemma}{original lemma}{original lemmas}
\Crefname{originallemma}{Original lemma}{Original lemmas}
\makeatletter
\newlabel{limits-proposition-characterize-locally-finite-presentation}{{limits-proposition-characterize-locally-finite-presentation}{source}{}{sourcepointer-5}{}}
\newlabel{limits-proposition-characterize-locally-finite-presentation@cref}{{[originallemma][0][]limits-proposition-characterize-locally-finite-presentation}{[1][0][]source}{}{}{}}
\makeatother
\crefname{originallemma}{original lemma}{original lemmas}
\Crefname{originallemma}{Original lemma}{Original lemmas}
\makeatletter
\newlabel{more-algebra-proposition-Noetherian-complete-G-ring}{{more-algebra-proposition-Noetherian-complete-G-ring}{source}{}{sourcepointer-6}{}}
\newlabel{more-algebra-proposition-Noetherian-complete-G-ring@cref}{{[originallemma][0][]more-algebra-proposition-Noetherian-complete-G-ring}{[1][0][]source}{}{}{}}
\makeatother
\crefname{originallemma}{original lemma}{original lemmas}
\Crefname{originallemma}{Original lemma}{Original lemmas}
\makeatletter
\newlabel{smoothing-theorem-popescu}{{smoothing-theorem-popescu}{source}{}{sourcepointer-7}{}}
\newlabel{smoothing-theorem-popescu@cref}{{[originallemma][0][]smoothing-theorem-popescu}{[1][0][]source}{}{}{}}
\makeatother
\title{Relative approximation of formal local maps before etale localisation}
\author{The Stacks Project Authors}
\date{stacks-project:more-morphisms.tex:local-snapshot-20261011; GNU FDL 1.2 or later}
\begin{document}
\maketitle
\noindent\textit{Editorial note.} Statement span is lines 11143-11203 of sources/stacks-more-morphisms/pinned-original-source.tex; proof span is lines 11205-11283, the complete author proof (77 body lines).\par\medskip
\noindent\textit{Editorial note.} Navigation stubs: 7 cited-but-not-included labels are supplied as non-load-bearing pointers with their pinned tags where the snapshot tag index has them.\par\medskip
\noindent\textit{Editorial note.} extra\_packages is empty on purpose: the retained preamble already loads xy with options and a second declaration clashes in the assembler wrapper.\par\medskip
\noindent\textit{Editorial note.} Text is retained verbatim, including the author's own typos and placeholders.\par\medskip
\noindent\textit{Allowed standard prerequisites (editorial).}\begin{itemize}
\item Explicit prerequisite boundary. Local maps of completions, finite type modifications, and the control lemmas for approximation; the reduction is retained.
\end{itemize}\medskip
\ifdefined\hypertarget\hypertarget{sourcepointer-1}{}\fi
\noindent\textit{Source reference.} Original lemma lemma-isomorphism-approximation: sources/stacks-more-morphisms/pinned-original-source.tex, source lines 11071--11095. Cited by the retained statement or proof but not part of the collected spans; supplied only as a navigation pointer. The argument does not depend on its content.\par\medskip
\ifdefined\hypertarget\hypertarget{sourcepointer-2}{}\fi
\noindent\textit{Source reference.} Original lemma lemma-slice-smooth: sources/stacks-more-morphisms/pinned-original-source.tex, source lines 10775--10783. Cited by the retained statement or proof but not part of the collected spans; supplied only as a navigation pointer. The argument does not depend on its content.\par\medskip
\ifdefined\hypertarget\hypertarget{sourcepointer-3}{}\fi
\noindent\textit{Source reference.} Original lemma limits-lemma-descend-etale: sources/stacks-more-morphisms/pinned-original-source.tex, source lines 11283--11283. Cited by the retained statement or proof but not part of the collected spans; supplied only as a navigation pointer. The argument does not depend on its content.\par\medskip
\ifdefined\hypertarget\hypertarget{sourcepointer-4}{}\fi
\noindent\textit{Source reference.} Original lemma limits-lemma-descend-finite-presentation: sources/stacks-more-morphisms/pinned-original-source.tex, source lines 11283--11283. Cited by the retained statement or proof but not part of the collected spans; supplied only as a navigation pointer. The argument does not depend on its content.\par\medskip
\ifdefined\hypertarget\hypertarget{sourcepointer-5}{}\fi
\noindent\textit{Source reference.} Original lemma limits-proposition-characterize-locally-finite-presentation: sources/stacks-more-morphisms/pinned-original-source.tex, source lines 11283--11283. Cited by the retained statement or proof but not part of the collected spans; supplied only as a navigation pointer. The argument does not depend on its content.\par\medskip
\ifdefined\hypertarget\hypertarget{sourcepointer-6}{}\fi
\noindent\textit{Source reference.} Original lemma more-algebra-proposition-Noetherian-complete-G-ring: sources/stacks-more-morphisms/pinned-original-source.tex, source lines 11283--11283. Cited by the retained statement or proof but not part of the collected spans; supplied only as a navigation pointer. The argument does not depend on its content.\par\medskip
\ifdefined\hypertarget\hypertarget{sourcepointer-7}{}\fi
\noindent\textit{Source reference.} Original lemma smoothing-theorem-popescu: sources/stacks-more-morphisms/pinned-original-source.tex, source lines 11283--11283. Cited by the retained statement or proof but not part of the collected spans; supplied only as a navigation pointer. The argument does not depend on its content.\par\medskip
\begin{lemma}
\label{lemma-relative-map-approximation-pre}
Let $X \to S$, $Y \to T$, $x$, $s$, $y$, $t$, $\sigma$, $y_\sigma$, and
$\varphi$ be given as follows: we have morphisms of schemes
$$
\vcenter{
\xymatrix{
X \ar[d] & Y \ar[d] \\
S & T
}
}
\quad\text{with points}\quad
\vcenter{
\xymatrix{
x \ar[d] & y \ar[d] \\
s & t
}
}
$$
Here $S$ is locally Noetherian and $T$ is of finite type over $\mathbf{Z}$.
The morphisms $X \to S$ and $Y \to T$ are locally of finite type.
The local ring $\mathcal{O}_{S, s}$ is a G-ring. The map
$$
\sigma : \mathcal{O}_{T, t} \longrightarrow \mathcal{O}_{S, s}^\wedge
$$
is a local homomorphism. Set
$Y_\sigma = Y \times_{T, \sigma} \Spec(\mathcal{O}_{S, s}^\wedge)$.
Next, $y_\sigma$ is a point of $Y_\sigma$ mapping to $y$ and
the closed point of $\Spec(\mathcal{O}_{S, s}^\wedge)$. Finally
$$
\varphi :
\mathcal{O}_{X, x}^\wedge
\longrightarrow
\mathcal{O}_{Y_\sigma, y_\sigma}^\wedge
$$
is an isomorphism of $\mathcal{O}_{S, s}^\wedge$-algebras.
In this situation there exists a commutative diagram
$$
\xymatrix{
X \ar[d] &
W \ar[l] \ar[rd] \ar[rr] & &
Y \times_{T, \tau} V \ar[r] \ar[ld] & Y \ar[d] \\
S & &
V \ar[ll] \ar[rr]^\tau & &
T
}
$$
of schemes and points $w \in W$, $v \in V$ such that
\begin{enumerate}
\item $(V, v) \to (S, s)$ is an elementary \'etale neighbourhood,
\item $(W, w) \to (X, x)$ is an elementary \'etale neighbourhood, and
\item $\tau(v) = t$.
\end{enumerate}
Let $y_\tau \in Y \times_T V$ correspond to $y_\sigma$
via the identification $(Y_\sigma)_s = (Y \times_T V)_v$.
Then
\begin{enumerate}
\item[(4)] $(W, w) \to (Y \times_{T, \tau} V, y_\tau)$ is an elementary
\'etale neighbourhood.
\end{enumerate}
\end{lemma}

\begin{proof}
Denote $X_\sigma = X \times_S \Spec(\mathcal{O}_{S, s}^\wedge)$
and $x_\sigma \in X_\sigma$ the unique point lying over $x$.
Observe that $\mathcal{O}_{S, s}^\wedge$ is a G-ring by
More on Algebra, Proposition
\ref{more-algebra-proposition-Noetherian-complete-G-ring}.
By Lemma \ref{lemma-isomorphism-approximation}
we can choose
$$
(X_\sigma, x_\sigma) \leftarrow (U, u) \rightarrow (Y_\sigma, y_\sigma)
$$
where both arrows are elementary \'etale neighbourhoods.

\medskip\noindent
After replacing $S$ by an open neighbourhood of $s$, we may
assume $S = \Spec(R)$ is affine. Since $\mathcal{O}_{S, s}$
is a G-ring by Smoothing Ring Maps, Theorem \ref{smoothing-theorem-popescu}
the ring $\mathcal{O}_{S, s}^\wedge$ is a filtered colimit of smooth
$R$-algebras. Thus we can write
$$
\Spec(\mathcal{O}_{S, s}^\wedge) = \lim S_i
$$
as a directed limit of affine schemes $S_i$ smooth over $S$.
Denote $s_i \in S_i$ the image of the closed point of
$\Spec(\mathcal{O}_{S, s}^\wedge)$. Observe that $\kappa(s) = \kappa(s_i)$.
Set $X_i = X \times_S S_i$ and denote $x_i \in X_i$ the unique
point mapping to $x$. Note that $\kappa(x) = \kappa(x_i)$.
Since $T$ is of finite type over $\mathbf{Z}$ by Limits, Proposition
\ref{limits-proposition-characterize-locally-finite-presentation}
we can choose an $i$ and a morphism $\sigma_i : (S_i, s_i) \to (T, t)$
of pointed schemes whose composition with
$\Spec(\mathcal{O}_{S, s}^\wedge) \to S_i$ is equal to $\sigma$.
Set $Y_i = Y \times_T S_i$ and denote $y_i$ the
image of $y_\sigma$. Note that $\kappa(y_i) = \kappa(y_\sigma)$.
By Limits, Lemma \ref{limits-lemma-descend-finite-presentation}
we can choose an $i$ and a diagram
$$
\xymatrix{
X_i \ar[rd] &
U_i \ar[l] \ar[d] \ar[r] &
Y_i \ar[ld] \\
& S_i
}
$$
whose base change to $\Spec(\mathcal{O}_{S, s}^\wedge)$
recovers $X_\sigma \leftarrow U \rightarrow Y_\sigma$.
By Limits, Lemma \ref{limits-lemma-descend-etale}
after increasing $i$ we may assume the morphisms
$X_i \leftarrow U_i \rightarrow Y_i$ are \'etale.
Let $u_i \in U_i$ be the image of $u$. Then $u_i \mapsto x_i$
hence
$\kappa(x) = \kappa(x_\sigma) = \kappa(u) \supset \kappa(u_i) \supset
\kappa(x_i) = \kappa(x)$ and we see that $\kappa(u_i) = \kappa(x_i)$.
Hence $(X_i, x_i) \leftarrow (U_i, u_i)$ is an elementary
\'etale neighbourhood. Since also $\kappa(y_i) = \kappa(y_\sigma) = \kappa(u)$
we see that also $(U_i, u_i) \to (Y_i, y_i)$ is an elementary
\'etale neighbourhood.

\medskip\noindent
At this point we have constructed a diagram
$$
\xymatrix{
X \ar[d] &
X \times_S S_i \ar[l] \ar[rd] &
U_i \ar[l] \ar[r] \ar[d] &
Y \times_T S_i \ar[r] \ar[ld] &
Y \ar[d] \\
S & &
S_i \ar[ll] \ar[rr] & &
T
}
$$
as in the statement of the lemma, except that $S_i \to S$ is smooth.
By Lemma \ref{lemma-slice-smooth} and after shrinking $S_i$
we can assume there exists a closed subscheme $V \subset S_i$ passing
through $s_i$ such that $V \to S$ is \'etale.
Setting $W$ equal to the scheme theoretic inverse image of $V$
in $U_i$ we conclude.
\end{proof}

\end{document}
